\chapter{Mengenbasierte Konstruktionen}
\label{b06guide:ch:constructions}

\section{Elementare Injektionen}

\begin{theorem}[Inklusion, Einermengen und Markierungen]
\label{b06guide:thm:elementary}
Für $A\subseteq B$ ist die Inklusion
\[
 \iota_{A,B}\colon A\rightarrowtail B,\qquad x\mapsto x
\]
injektiv. Insbesondere ist $\operatorname{id}_A$ injektiv.
Für jede Menge $A$ und jedes fest gewählte Objekt $c$
sind außerdem folgende Funktionen injektiv:
\[
 \begin{aligned}
 A&\longrightarrow\mathcal P(A),&x&\longmapsto\{x\},\\
 A&\longrightarrow A\times\{c\},&x&\longmapsto(x,c),\\
 A&\longrightarrow A\times A,&x&\longmapsto(x,x).
 \end{aligned}
\]
\end{theorem}

\begin{theorem}[Graphabbildung]
\label{b06guide:thm:graphmap}
Für jede Funktion $h\colon A\to B$ ist die Funktion
\[
 \gamma_h\colon A\rightarrowtail A\times B,\qquad
 x\mapsto(x,h(x))
\]
injektiv. Es wird keine Injektivität von $h$ vorausgesetzt.
Die Aussage liefert eine Injektion aus einer beliebigen Funktion.
\end{theorem}

\section{Produkte und gemeinsame Unterscheidung}

\begin{theorem}[Produkt zweier Injektionen]
\label{b06guide:thm:product}
Seien $f\colon A\rightarrowtail B$ und $g\colon C\rightarrowtail D$.
Dann ist
\[
 f\times g\colon A\times C\rightarrowtail B\times D,\qquad
 (a,c)\mapsto(f(a),g(c))
\]
injektiv.

Für beliebige Funktionen $f\colon A\to B$ und $g\colon C\to D$
gilt bei $A\ne\varnothing$ und $C\ne\varnothing$ auch die Umkehrung:
Die Produktabbildung ist genau dann injektiv, wenn beide Faktoren
injektiv sind. Ist $A$ oder $C$ leer, ist die Produktabbildung
stets injektiv.
\end{theorem}

\begin{theorem}[Gemeinsam unterscheidende Komponenten]
\label{b06guide:thm:pairing}
Für $f\colon A\to B$ und $g\colon A\to C$ ist
\[
 h\colon A\to B\times C,\qquad x\mapsto(f(x),g(x))
\]
genau dann injektiv, wenn
\[
 \forall x,y\in A\,
 \bigl(f(x)=f(y)\land g(x)=g(y)\Rightarrow x=y\bigr).
\]
Insbesondere genügt die Injektivität einer der beiden Komponenten.
\end{theorem}

\begin{definition}[Indiziertes Produkt und markierte Vereinigung]
Für eine gegebene Familie $(A_i)_{i\in I}$ seien
\[
 \begin{aligned}
 \prod_{i\in I}A_i
   &=\left\{a\colon I\to\bigcup_{i\in I}A_i\
                 \middle|\ \forall i\in I\ (a(i)\in A_i)\right\},\\
 \coprod_{i\in I}A_i
   &=\{(i,x)\mid i\in I,\ x\in A_i\}.
 \end{aligned}
\]
Für $I=\varnothing$ ist das Produkt $\{\varnothing\}$
und die markierte Vereinigung leer.
\end{definition}

\begin{theorem}[Indizierte Produkte gegebener Injektionen]
\label{b06guide:thm:indexedproduct}
Sei $(f_i\colon A_i\rightarrowtail B_i)_{i\in I}$ eine
\emph{gegebene} Familie von Injektionen. Dann ist
\[
 F\colon\prod_{i\in I}A_i\rightarrowtail\prod_{i\in I}B_i,
 \qquad F(a)(i)=f_i(a(i))
\]
injektiv. Die Aussage gilt auch für eine leere Indexmenge
und für ein leeres Ausgangsprodukt; sie benötigt keine
Nichtleerheitsbehauptung über unendliche Produkte.

Für eine gegebene Familie beliebiger Funktionen
$f_i\colon A_i\to B_i$ gilt bei
$\prod_{i\in I}A_i\ne\varnothing$ auch die Umkehrung:
Ist $F$ injektiv, so ist jedes $f_i$ injektiv.
\end{theorem}

\begin{theorem}[Markierte Vereinigungen gegebener Injektionen]
\label{b06guide:thm:taggedunion}
Sei $(f_i\colon A_i\rightarrowtail B_i)_{i\in I}$ eine
gegebene Familie von Injektionen. Dann ist
\[
 F\colon\coprod_{i\in I}A_i\rightarrowtail\coprod_{i\in I}B_i,
 \qquad (i,x)\mapsto(i,f_i(x))
\]
injektiv. Die Mengen $A_i$ untereinander und die Mengen $B_i$
untereinander müssen nicht disjunkt sein.
Für eine gegebene Familie beliebiger Funktionen $f_i$
ist $F$ genau dann injektiv, wenn jedes $f_i$ injektiv ist.
\end{theorem}

\section{Funktionenmengen und Potenzmengen}

\begin{theorem}[Nachkomposition auf einer Funktionenmenge]
\label{b06guide:thm:functionsets}
Für $f\colon A\rightarrowtail B$ und eine beliebige Menge $T$
ist die Funktion
\[
 f^T\colon A^T\rightarrowtail B^T,\qquad u\mapsto f\circ u
\]
injektiv. Für $T\ne\varnothing$ gilt für beliebiges $f\colon A\to B$
auch die Umkehrung: Ist $f^T$ injektiv, so ist $f$ injektiv.
Für $T=\varnothing$ ist $f^T$ stets injektiv.
\end{theorem}

\begin{theorem}[Bildfunktion als neue Injektion]
\label{b06guide:thm:powersetinjection}
Jede Injektion $f\colon A\rightarrowtail B$ liefert die Injektion
\[
 f_*\colon\mathcal P(A)\rightarrowtail\mathcal P(B),\qquad
 U\mapsto f[U].
\]
Diese Konstruktion darf erneut angewandt werden. Insbesondere ist
\[
 \mathcal P(\mathcal P(A))
   \longrightarrow\mathcal P(\mathcal P(B)),\qquad
 \mathcal U\longmapsto\{f[U]\mid U\in\mathcal U\}
\]
injektiv.
\end{theorem}

\begin{theorem}[Adjunktion einer disjunkten festen Menge]
\label{b06guide:thm:adjunction}
Seien $A,C$ Mengen mit $A\cap C=\varnothing$.
Dann ist
\[
 J_C\colon\mathcal P(A)\rightarrowtail\mathcal P(A\cup C),
 \qquad U\mapsto U\cup C
\]
injektiv. Ohne Voraussetzung an $A,C$ ist diese Abbildung
genau dann injektiv, wenn $A\cap C=\varnothing$.
Die Aussage betrifft nur Mengen; eine Abgeschlossenheitsstruktur
auf einer Mengenfamilie wird nicht vorausgesetzt.
\end{theorem}